\appendix

\chapter{Preliminary Development Sample}
\label{app:preliminary-sample}

This appendix preserves the operational record that was accumulated while the
prototype was being built, between 28~June and 9~July~2026. It comprises 17
invocations of the hybrid pipeline orchestrator
(\texttt{scripts/run\_hybrid\_pipeline.py}) together with their persisted gate,
approval, and rejection records, retained in \texttt{logs/}.

This material is reported separately from the results in
Chapter~\ref{Chapter4} and supports no claim made there. The reason is
methodological rather than a matter of data quality. These runs were a
convenience sample: they were produced while features were being developed and
debugged, so they over-represent whatever happened to be under work at the time;
several exercised projects that were incomplete; and the decision as to which
invocations to count would necessarily have been made after seeing their
outcomes. A sample selected in that way cannot distinguish a system that behaves
correctly from an analyst who selected the runs in which it did. The
pre-registered campaign described in
Section~\ref{subsec:evaluation-campaign} was constructed to replace it.

The sample is retained for two reasons. It documents the development history
honestly, showing what the evidence base looked like before the campaign
methodology was adopted. It also contains something the registered campaign does
not: real deployments against a live on-premises host, which the offline
campaign deliberately does not attempt.

\section{Observed Counts}

Of the 17 recorded invocations, 16 produced a persisted run summary. The
remaining invocation wrote only a log header before terminating without a
summary; the orchestrator's exception path was not captured at that stage of
development, so no cause is retained for it. Each completed invocation
independently gates a CDK branch, an Ansible branch, or both, yielding 24
branch-level gate decisions --- 9 CDK and 15 Ansible --- across the 16 completed
runs.

\begin{table}[ht]
\centering
\caption{Counts observed in the preliminary development sample.}
\label{tab:app-preliminary-counts}
\begin{tabular}{|L{5.9cm}|C{3.0cm}|L{4.7cm}|}
\hline
\textbf{Measure} & \textbf{Result} & \textbf{Evidence source} \\
\hline
Total pipeline runs & $17$ & Execution records \\
\hline
Runs reaching a valid terminal decision & $16$ & Gate reports \\
\hline
End-to-end completion rate & $94.1\%$ ($16/17$) & Calculated from run records \\
\hline
Pass, review, and reject outcomes (branch-level, $n=24$) & $18$ / $4$ / $2$ &
Persisted decision records \\
\hline
Unexpected pipeline failures & $2$ & Stage results and execution logs \\
\hline
Rejected changes incorrectly deployed & $0$ of $2$ & Gate and deployment records \\
\hline
Review cases executed without approval & $0$ of $4$ & Approval and deployment records \\
\hline
Runs continuing after a cost-scanner error & $5$ of $5$ & Scanner status and gate reports \\
\hline
\end{tabular}
\end{table}

Two unexpected failures were observed. The first is the incomplete invocation
described above. The second occurred within a completed run: an Ansible branch
passed its security gate with a score of $0$, but the subsequent deployment step
returned a non-zero exit code, giving the terminal status
\texttt{deploy\_failed}. Neither resulted from a security-gate rejection, which
is why both are counted as failures rather than as correct refusals.

Both reject decisions were withheld from deployment and each has a matching
rejection record. Of the four review decisions, three remained unapproved and
undeployed and one was deployed after an explicit approval record was written.
Five CDK branches encountered an error from the Infracost cost-analysis scanner
--- a timeout or an unexpected response schema --- and all five still produced a
valid decision, which was the first indication that the gate tolerates an
unavailable optional component rather than terminating.

\section{Live Deployment Evidence}

Fifteen Ansible-branch invocations appear in this sample, of which seven
proceeded to an actual deployment attempt against the on-premises virtual
machine, addressed by its Tailscale mesh address in the checked-in inventory
(\texttt{examples/hybrid-webapp-demo/ansible/inventory.ini}). Six of the seven
completed successfully. One returned a non-zero exit code whose specific cause
--- connectivity, authentication, or a task-level condition --- cannot be
distinguished from the retained artifacts.

These seven attempts are the only evidence in the project of the apply
checkpoint being reached against a live target. They are too few, and too
loosely controlled, to support an attainment rate, and post-apply state
verification and repeated-run idempotency were not captured at all. The three
checkpoints reported as evidence gaps in
Section~\ref{subsec:results-rq2} therefore remain gaps: this sample narrows them
somewhat but does not close them.

\section{What the Sample Could Not Show}

Every Ansible-branch decision in this sample was a \emph{pass} with a score of
zero. No review or reject decision was ever observed on the on-premises path, so
the sample could not show that the gate constrains that path in the same way it
constrains the cloud path; that property had to be inferred from the shared
decision function rather than observed.

This limitation was the direct motivation for the calibrated fixtures described
in Section~\ref{subsec:evaluation-campaign}. An evaluation corpus in which the
subject never encounters a case it should refuse cannot demonstrate that it
refuses correctly. The registered campaign authored fixtures for each band on
both branches precisely so that the refusal paths would be exercised rather than
assumed.
